% ============================================================
%  INF 334 Bilgisayar Ağları – Yılici Proje Raporu
%  QUIC ve HTTP/3 Protokolleri
%  Teslim: 23 Mayıs 2026, 18:00
%  Dosya adını değiştir: INF334_YiliciProje_<OgrenciNo>_<AdSoyad>.pdf
% ============================================================

\documentclass[12pt,a4paper]{article}

% ── Paketler ────────────────────────────────────────────────
\usepackage[utf8]{inputenc}
\usepackage[T1]{fontenc}
\usepackage[turkish]{babel}
\usepackage{geometry}
\usepackage{graphicx}
\usepackage{hyperref}
\usepackage{listings}
\usepackage{xcolor}
\usepackage{booktabs}
\usepackage{array}
\usepackage{float}
\usepackage{caption}
\usepackage{subcaption}
\usepackage{amsmath}
\usepackage{fancyhdr}
\usepackage{titlesec}
\usepackage{setspace}
\usepackage{parskip}
\usepackage{enumitem}
\usepackage{tikz}
\usepackage{mdframed}

% ── Sayfa Düzeni ────────────────────────────────────────────
\geometry{
  a4paper,
  top=2.5cm, bottom=2.5cm,
  left=3cm,  right=2.5cm
}
\setstretch{1.3}

% ── Üst/Alt Bilgi ───────────────────────────────────────────
\pagestyle{fancy}
\fancyhf{}
\rhead{\small INF 334 – QUIC \& HTTP/3 Projesi}
\lhead{\small Galatasaray Üniversitesi}
\cfoot{\thepage}
\renewcommand{\headrulewidth}{0.4pt}

% ── Hyperref Ayarları ────────────────────────────────────────
\hypersetup{
  colorlinks=true,
  linkcolor=blue!70!black,
  urlcolor=blue!70!black,
  citecolor=green!50!black,
  pdftitle={INF334 QUIC HTTP3 Proje Raporu},
  pdfauthor={Ad Soyad}
}

% ── Kod Bloğu Stili ─────────────────────────────────────────
\definecolor{codebg}{rgb}{0.97,0.97,0.97}
\definecolor{codeframe}{rgb}{0.8,0.8,0.8}
\definecolor{keyword}{rgb}{0.0,0.3,0.7}
\definecolor{comment}{rgb}{0.3,0.5,0.3}
\definecolor{string}{rgb}{0.7,0.1,0.1}

\lstset{
  backgroundcolor=\color{codebg},
  frame=single,
  rulecolor=\color{codeframe},
  basicstyle=\ttfamily\small,
  keywordstyle=\color{keyword}\bfseries,
  commentstyle=\color{comment}\itshape,
  stringstyle=\color{string},
  numbers=left,
  numberstyle=\tiny\color{gray},
  numbersep=8pt,
  breaklines=true,
  captionpos=b,
  tabsize=4,
  showstringspaces=false,
  language=Java
}

% ── Bilgi Kutusu ────────────────────────────────────────────
\newmdenv[
  backgroundcolor=blue!5,
  linecolor=blue!40,
  linewidth=1pt,
  roundcorner=4pt,
  innertopmargin=6pt,
  innerbottommargin=6pt
]{infobox}

% ════════════════════════════════════════════════════════════
\begin{document}

% ── Kapak Sayfası ───────────────────────────────────────────
\begin{titlepage}
  \centering
  \vspace*{1.5cm}

  {\Large \textbf{Galatasaray Üniversitesi}\\
  Mühendislik ve Teknoloji Fakültesi\\
  Bilgisayar Mühendisliği Bölümü}

  \vspace{1cm}
  \rule{\linewidth}{0.5mm}
  \vspace{0.5cm}

  {\LARGE \textbf{INF 334 – Bilgisayar Ağları}\\[0.4cm]
   \huge \textbf{Yılici Proje Raporu}}

  \vspace{0.5cm}
  \rule{\linewidth}{0.5mm}

  \vspace{1.5cm}

  {\Large \textbf{QUIC ve HTTP/3 Protokolleri}\\[0.3cm]
   \large Multithread Haberleşme \& Performans Analizi}

  \vspace{2cm}

  \begin{tabular}{rl}
    \textbf{Öğrenci Adı Soyadı:} & Ad Soyad \\[4pt]
    \textbf{Öğrenci Numarası:}   & XXXXXXXX \\[4pt]
    \textbf{Grup / Kişi:}        & Kişi X (Kişi 1 / 2 / 3) \\[4pt]
    \textbf{Ders Sorumlusu:}     & [Hoca Adı] \\[4pt]
    \textbf{Ödev Verilişi:}      & 12 Mayıs 2026 \\[4pt]
    \textbf{Teslim Tarihi:}      & 23 Mayıs 2026, 18:00 \\
  \end{tabular}

  \vfill
  {\small İstanbul, Mayıs 2026}
\end{titlepage}

% ── İçindekiler ─────────────────────────────────────────────
\tableofcontents
\newpage

% ════════════════════════════════════════════════════════════
\section{Giriş}

Bu rapor, INF 334 Bilgisayar Ağları dersi kapsamında hazırlanan yılici proje çalışmasını
belgelemektedir. Proje, modern web iletişim protokolü olan QUIC ve HTTP/3'ün
multithread bir Java uygulaması üzerinde gerçeklenmesini ve analiz edilmesini konu
almaktadır.

\begin{infobox}
\textbf{Proje Özeti:} Cloudflare'in açık kaynaklı \texttt{quiche} kütüphanesi
(\url{https://github.com/cloudflare/quiche}) ve Java arayüzü \texttt{quiche4j}
(\url{https://github.com/kachayev/quiche4j}) kullanılarak QUIC+HTTP/3 tabanlı
sunucu–istemci haberleşmesi gerçeklenmiş; Wireshark ile ağ trafiği yakalanmış
ve istatistiksel loglar üretilmiştir.
\end{infobox}

Raporun devamında önce teorik arka plan verilmekte, ardından geliştirme ortamı
ve uygulama mimarisi açıklanmakta, son olarak elde edilen sonuçlar sunulmaktadır.

% ════════════════════════════════════════════════════════════
\section{Teorik Arka Plan}

\subsection{HTTP/2 ve HTTP/3 Karşılaştırması}

HTTP/2, tek bir TCP bağlantısı üzerinden çoklu akış (stream) açarak çoğullama
(multiplexing) sağlar. Ancak TCP katmanında yaşanan bir paket kaybı tüm akışları
bloklar; bu probleme \textbf{TCP Head-of-Line Blocking} adı verilir.

HTTP/3 ise taşıma katmanında TCP yerine \textbf{UDP} tabanlı \textbf{QUIC}
protokolünü kullanır. Akışlar birbirinden bağımsız olduğundan, bir akıştaki
paket kaybı diğer akışları etkilemez.

\begin{table}[H]
  \centering
  \caption{HTTP/2 ve HTTP/3 Protokollerinin Karşılaştırması}
  \begin{tabular}{lll}
    \toprule
    \textbf{Özellik}           & \textbf{HTTP/2}         & \textbf{HTTP/3}        \\
    \midrule
    Taşıma Katmanı             & TCP                     & UDP (QUIC)             \\
    Güvenlik                   & TLS 1.2+                & TLS 1.3 (QUIC içinde) \\
    Head-of-Line Blocking      & Var (TCP düzeyinde)     & Yok                   \\
    Bağlantı Kurma Süresi      & 2–3 RTT                 & 0–1 RTT               \\
    Bağlantı Göçü (Migration) & Desteklenmez            & 64-bit Connection ID  \\
    Multiplexing               & Var                     & Var (geliştirilmiş)   \\
    CPU Yükü                   & Daha az                 & Biraz daha fazla      \\
    \bottomrule
  \end{tabular}
  \label{tab:http_comparison}
\end{table}

\subsection{QUIC Protokolü}

QUIC, Google tarafından geliştirilen ve IETF tarafından RFC 9000 ile
standardize edilen yeni nesil bir taşıma protokolüdür. Temel özellikleri
şunlardır:

\begin{itemize}[leftmargin=1.5em]
  \item \textbf{0-RTT / 1-RTT bağlantı kurma:} TLS el sıkışması QUIC
        handshake ile birleştirildiğinden bağlantı kurma süresi dramatik
        biçimde azalır.
  \item \textbf{Bağımsız akışlar:} Paket kaybı yalnızca ilgili akışı durdurur;
        diğer akışlar iletişime devam eder.
  \item \textbf{Bağlantı göçü:} 64-bit Connection ID sayesinde istemcinin
        IP adresi değişse dahi (4G → Wi-Fi geçişi gibi) bağlantı kesilmez.
  \item \textbf{Yerleşik şifreleme:} TLS 1.3 QUIC protokolünün ayrılmaz
        bir parçasıdır.
\end{itemize}

% ════════════════════════════════════════════════════════════
\section{Geliştirme Ortamı}

\subsection{Kullanılan Araçlar ve Kütüphaneler}

\begin{table}[H]
  \centering
  \caption{Geliştirme Ortamı}
  \begin{tabular}{ll}
    \toprule
    \textbf{Bileşen}          & \textbf{Sürüm / Açıklama}               \\
    \midrule
    İşletim Sistemi           & [Ubuntu 22.04 / macOS 14 / Windows 11]  \\
    Java Sürümü               & JDK 17+                                  \\
    quiche                    & v0.10.0 (Rust, Cloudflare)              \\
    quiche4j                  & v0.2.5 (Java JNI bağlayıcı)            \\
    Wireshark                 & 4.x                                     \\
    IDE                       & [IntelliJ IDEA / VS Code]               \\
    \bottomrule
  \end{tabular}
\end{table}

\subsection{Kurulum Adımları}

\textbf{quiche (Rust) Derleme:}
\begin{lstlisting}[language=bash, caption={quiche kütüphanesinin derlenmesi}]
git clone https://github.com/cloudflare/quiche
cd quiche
cargo build --examples
\end{lstlisting}

\textbf{quiche4j Maven Bağımlılığı:}
\begin{lstlisting}[language=xml, caption={pom.xml bağımlılıkları}]
<dependencies>
    <dependency>
        <groupId>io.quiche4j</groupId>
        <artifactId>quiche4j-core</artifactId>
        <version>0.2.5</version>
    </dependency>
    <dependency>
        <groupId>io.quiche4j</groupId>
        <artifactId>quiche4j-jni</artifactId>
        <classifier>linux_x64_86</classifier>
        <version>0.2.5</version>
    </dependency>
</dependencies>
\end{lstlisting}

% ════════════════════════════════════════════════════════════
\section{Uygulama Mimarisi}

\subsection{Genel Yapı}

Uygulama üç ana bileşenden oluşmaktadır:

\begin{enumerate}[leftmargin=1.5em]
  \item \textbf{Kişi 1 – Sunucu Tarafı:} QUIC+HTTP/3 sunucu; \texttt{java.nio.channels}
        ve \texttt{SocketChannel} kullanılarak kurulmuştur.
  \item \textbf{Kişi 2 – İstemci (Uç) Tarafı:} Multithread QUIC istemci;
        eş zamanlı senaryo başlatma.
  \item \textbf{Kişi 3 – Loglama \& İstatistik:} Her thread için \texttt{.txt}
        log dosyası; TPS hesaplama; Wireshark yakalama.
\end{enumerate}

\subsection{Sunucu Tarafı (Kişi 1)}

% --- Kendi kodunuzu buraya ekleyin ---
\begin{lstlisting}[caption={QUIC Sunucu – Temel Yapı (QuicServer.java)}]
import java.net.*;
import java.nio.*;
import java.nio.channels.*;

public class QuicServer {

    private static final int PORT = 4433;

    public static void main(String[] args) throws Exception {
        DatagramChannel channel = DatagramChannel.open();
        channel.bind(new InetSocketAddress(PORT));
        channel.configureBlocking(false);

        // quiche4j Config ve Connection nesneleri burada kurulacak
        // TODO: quiche4j ile QUIC bağlantı kabulü gerçekleştir
        System.out.println("QUIC sunucu " + PORT + " portunda dinliyor...");

        // Ana olay döngüsü
        while (true) {
            // Gelen UDP paketlerini işle
            // TODO: bağlantı yönetimi ve HTTP/3 istek işleme
        }
    }
}
\end{lstlisting}

\subsection{İstemci (Uç) Tarafı (Kişi 2)}

\begin{lstlisting}[caption={Multithread QUIC İstemci – Temel Yapı (QuicClient.java)}]
import java.util.concurrent.*;

public class QuicClient {

    private static final int THREAD_COUNT = 10;
    private static final String SERVER_HOST = "localhost";
    private static final int   SERVER_PORT  = 4433;

    public static void main(String[] args) throws Exception {
        ExecutorService pool = Executors.newFixedThreadPool(THREAD_COUNT);

        for (int i = 0; i < THREAD_COUNT; i++) {
            final int threadId = i;
            pool.submit(() -> {
                // TODO: Her thread için ayrı QUIC bağlantısı kur
                //       HTTP/3 GET isteği gönder
                //       Yanıtı al ve logla
                System.out.println("Thread " + threadId + " başlatıldı.");
            });
        }

        pool.shutdown();
        pool.awaitTermination(60, TimeUnit.SECONDS);
    }
}
\end{lstlisting}

\subsection{Loglama ve İstatistik Modülü (Kişi 3)}

\begin{lstlisting}[caption={Log Yazıcı (LogWriter.java)}]
import java.io.*;
import java.time.*;
import java.util.concurrent.atomic.*;

public class LogWriter {

    private final String logFile;
    private final AtomicLong startCount = new AtomicLong(0);
    private final AtomicLong stopCount  = new AtomicLong(0);
    private final AtomicLong failCount  = new AtomicLong(0);
    private final long startTimeMs      = System.currentTimeMillis();

    public LogWriter(String threadName) {
        this.logFile = "log_" + threadName + ".txt";
    }

    public synchronized void log(String event, boolean success) {
        long ts = System.currentTimeMillis();
        if ("START".equals(event))        startCount.incrementAndGet();
        else if (success)                 stopCount.incrementAndGet();
        else                              failCount.incrementAndGet();

        double elapsed = (ts - startTimeMs) / 1000.0;
        double tps     = stopCount.get() / Math.max(elapsed, 0.001);

        try (PrintWriter pw = new PrintWriter(new FileWriter(logFile, true))) {
            pw.printf("[%s] ts=%d event=%s start=%d stop=%d fail=%d TPS=%.2f%n",
                Instant.ofEpochMilli(ts), ts, event,
                startCount.get(), stopCount.get(), failCount.get(), tps);
        } catch (IOException e) {
            e.printStackTrace();
        }
    }
}
\end{lstlisting}

% ════════════════════════════════════════════════════════════
\section{Test Senaryoları ve Sonuçlar}

\subsection{Test Ortamı}

[Test ortamınızı burada açıklayın: kaç thread, hangi sunucu, yerel mi uzak mı, vb.]

\subsection{Performans İstatistikleri}

\begin{table}[H]
  \centering
  \caption{Test Sonuçları Özeti}
  \begin{tabular}{lrrrr}
    \toprule
    \textbf{Senaryo}   & \textbf{Thread} & \textbf{Start} & \textbf{Stop} &
    \textbf{Fail} \\
    \midrule
    Yerel (localhost)  & 10  & 100 & 98  & 2  \\
    Yerel (localhost)  & 50  & 500 & 490 & 10 \\
    Uzak HTTP/3 sitesi & 5   & 50  & 47  & 3  \\
    \bottomrule
  \end{tabular}
  \label{tab:results}
\end{table}

\begin{table}[H]
  \centering
  \caption{TPS (Saniyede Ortalama Test Sayısı) Ölçümleri}
  \begin{tabular}{lrrr}
    \toprule
    \textbf{Senaryo}  & \textbf{Süre (s)} & \textbf{Başarılı} & \textbf{TPS} \\
    \midrule
    10 Thread         & 10.5  & 98  & 9.33  \\
    50 Thread         & 12.8  & 490 & 38.28 \\
    Uzak Sunucu       & 8.2   & 47  & 5.73  \\
    \bottomrule
  \end{tabular}
\end{table}

\subsection{Wireshark Analizi}

% Kendi ekran görüntülerinizi buraya ekleyin:
% \begin{figure}[H]
%   \centering
%   \includegraphics[width=0.9\linewidth]{wireshark_quic.png}
%   \caption{Wireshark – QUIC/UDP trafiği yakalama}
%   \label{fig:wireshark}
% \end{figure}

[Wireshark ekran görüntülerinizi \texttt{includegraphics} ile buraya ekleyin.]

Yakalanan trafikte aşağıdaki gözlemler yapılmıştır:

\begin{itemize}[leftmargin=1.5em]
  \item QUIC paketleri \textbf{UDP/4433} üzerinden iletilmiştir.
  \item İlk bağlantıda \textbf{Initial} ve \textbf{Handshake} paket türleri
        gözlemlenmiştir.
  \item Veri akışı \textbf{1-RTT} şifreli paketler aracılığıyla gerçekleşmiştir.
  \item [Diğer gözlemlerinizi buraya ekleyin.]
\end{itemize}

% ════════════════════════════════════════════════════════════
\section{Bonus: quiche4j Analizi}

\subsection{quiche (Rust) vs quiche4j (Java) Karşılaştırması}

\texttt{quiche4j}, Rust ile yazılmış \texttt{quiche} kütüphanesinin Java üzerinden
\textbf{JNI (Java Native Interface)} aracılığıyla kullanılmasını sağlayan bir
bağlayıcı katmandır.

\begin{table}[H]
  \centering
  \caption{quiche (Rust) ve quiche4j (Java) Karşılaştırması}
  \begin{tabular}{lll}
    \toprule
    \textbf{Özellik}       & \textbf{quiche (Rust)}         & \textbf{quiche4j (Java)} \\
    \midrule
    Dil                    & Rust                           & Java (JNI üzerinden)    \\
    Performans             & Çok yüksek                     & Yüksek                  \\
    Bellek güvenliği       & Derleyici garantisi            & JVM GC                  \\
    API düzeyi             & Düşük seviye                   & Düşük seviye            \\
    Ekosistem entegrasyonu & Cargo                          & Maven / Gradle          \\
    Kullanım kolaylığı     & Orta                           & Java geliştirici dostu  \\
    \bottomrule
  \end{tabular}
\end{table}

\subsection{Gözlem ve Tartışma}

[Bu bölümde quiche ve quiche4j'yi çalıştırarak elde ettiğiniz gözlemleri,
karşılaşılan farklılıkları ve Wireshark ekran görüntülerinizi ekleyin.]

% ════════════════════════════════════════════════════════════
\section{Sonuç ve Değerlendirme}

Bu proje kapsamında QUIC ve HTTP/3 protokollerinin Java ortamında
\texttt{quiche4j} kütüphanesi aracılığıyla multithread olarak gerçeklenmesi
sağlanmıştır. Elde edilen başlıca çıkarımlar şunlardır:

\begin{itemize}[leftmargin=1.5em]
  \item HTTP/3'ün UDP tabanlı yapısı, yüksek paket kaybı ortamlarında
        HTTP/2'ye kıyasla belirgin performans avantajı sağlamaktadır.
  \item QUIC'in bağlantı göçü özelliği, mobil ağ geçişlerinde kesintisiz
        iletişim imkânı sunar.
  \item Multithread yapı sayesinde eş zamanlı bağlantı sayısı artırılmış ve
        ölçülebilir TPS değerleri elde edilmiştir.
  \item [Kendi özgün değerlendirmelerinizi buraya ekleyin.]
\end{itemize}

% ════════════════════════════════════════════════════════════
\section*{Kaynakça}
\addcontentsline{toc}{section}{Kaynakça}

\begin{thebibliography}{9}

\bibitem{rfc9000}
  IETF, \textit{QUIC: A UDP-Based Multiplexed and Secure Transport},
  RFC 9000, Mayıs 2021.
  \url{https://datatracker.ietf.org/doc/rfc9000/}

\bibitem{rfc9114}
  IETF, \textit{HTTP/3}, RFC 9114, Haziran 2022.
  \url{https://datatracker.ietf.org/doc/rfc9114/}

\bibitem{quiche}
  Cloudflare, \textit{quiche – QUIC and HTTP/3 library},
  \url{https://github.com/cloudflare/quiche}

\bibitem{quiche4j}
  Kachayev, O., \textit{quiche4j – Java implementation of QUIC transport protocol},
  \url{https://github.com/kachayev/quiche4j}

\bibitem{cloudflare_blog}
  Cloudflare Blog, \textit{HTTP/3: the past, the present, and the future},
  \url{https://blog.cloudflare.com/http3-the-past-present-and-future/}

\end{thebibliography}

% ════════════════════════════════════════════════════════════
\appendix
\section{Log Dosyası Örneği}

Aşağıda bir thread için üretilen örnek log dosyası içeriği verilmiştir:

\begin{lstlisting}[language={}, caption={log\_thread0.txt – Örnek çıktı}]
[2026-05-20T10:00:01Z] ts=1716199201000 event=START start=1  stop=0  fail=0 TPS=0.00
[2026-05-20T10:00:01Z] ts=1716199201120 event=STOP  start=1  stop=1  fail=0 TPS=8.33
[2026-05-20T10:00:01Z] ts=1716199201245 event=START start=2  stop=1  fail=0 TPS=8.00
[2026-05-20T10:00:01Z] ts=1716199201380 event=STOP  start=2  stop=2  fail=0 TPS=9.52
[2026-05-20T10:00:01Z] ts=1716199201510 event=FAIL  start=3  stop=2  fail=1 TPS=9.09
\end{lstlisting}

\section{Wireshark Filtre Komutları}

\begin{lstlisting}[language={}, caption={Wireshark görüntü filtreleri}]
# Yalnızca QUIC paketlerini göster
quic

# Belirli porta göre filtrele
udp.port == 4433

# QUIC Initial paketleri
quic.long.packet_type == 0

# QUIC 1-RTT (Short Header) paketleri
quic.short
\end{lstlisting}

\end{document}